% Specification review / feedback for SQIsign authors
% Derived from implementing sqisign-selkie (Rust, NIST-I)
\documentclass[a4paper,11pt]{article}

\usepackage{amsfonts,amsmath,amssymb}
\usepackage[T1]{fontenc}
\usepackage{libertinus}
\usepackage[libertine]{newtxmath}
\usepackage{microtype}
\usepackage[dvipsnames]{xcolor}
\usepackage{enumitem}
\usepackage{booktabs}
\usepackage[colorlinks=true, urlcolor=BrickRed, linkcolor=BrickRed,
            citecolor=BrickRed]{hyperref}
\usepackage{xurl}
\usepackage{geometry}
\geometry{margin=28mm, marginparwidth=0pt}

\newcommand{\algo}[1]{Algorithm~#1}
\newcommand{\Fp}{\mathbb{F}_p}
\newcommand{\Fpii}{\mathbb{F}_{p^2}}

% Severity badges — inline colored small-caps
\newcommand{\critical}{\textsc{\color{BrickRed}Critical}}
\newcommand{\high}{\textsc{\color{BurntOrange}High}}
\newcommand{\medium}{\textsc{\color{RoyalBlue}Medium}}
\newcommand{\low}{\textsc{\color{OliveGreen}Low}}

% Finding header: severity + affected algorithms on one line
\newcommand{\finding}[2]{\noindent\textit{#1}\hfill#2\par\smallskip}

\title{SQIsign v2 Specification Review\\[0.3em]
  {\large Feedback from an Independent Rust Implementation}}
\author{Deirdre Connolly\\[0.2em]
  \small Selkie Cryptography\\
  \small\texttt{durumcrustulum@gmail.com}}
\date{\today\enspace{\small(living document)}}

\begin{document}
\maketitle

\begin{abstract}
Each item below describes something that the SQIsign
specification~\cite{sqisign-spec} (version 2025-07-07) leaves unstated,
ambiguous, or implicit, and that we could only resolve by inspecting the
C reference implementation.  We file these as constructive suggestions so
that future implementors can achieve interoperability from the spec
alone.
\end{abstract}

\tableofcontents

% =====================================================================
\section{Severity scale}

\smallskip
\begin{tabular}{@{}lp{0.86\textwidth}@{}}
\toprule
\critical & The spec as written produces a non-interoperable
    implementation: verification rejects valid signatures, or signing
    produces unverifiable ones.  The spec must be amended. \\[0.4em]
\high & An implementor will almost certainly write a bug that passes
    unit tests but fails KAT or cross-implementation testing.
    Debugging requires reading the C reference source. \\[0.4em]
\medium & An implementor may write a bug depending on reasonable but
    wrong assumptions about conventions.  A short note eliminates the
    ambiguity. \\[0.4em]
\low & The spec is correct but incomplete.  A remark or forward pointer
    would help. \\
\bottomrule
\end{tabular}

% =====================================================================
\section{\texorpdfstring{$\Fpii$}{Fp2} square root canonicalization}
\label{sec:sqrt}

\finding{Affects \algo{8.12} (\textsc{DifferencePoint}), and
transitively every algorithm that consumes a torsion basis.}{\critical}

\algo{8.12} says ``compute $\delta = \sqrt{B_{XZ}^2 - B_{XX} B_{ZZ}}$''
without specifying which of the two square roots in $\Fpii$ to use.
The two roots $\delta$ and $-\delta$ produce two \emph{affinely distinct}
points $x_{P-Q}$.  These are not projectively equivalent; they are
genuinely different points.  The wrong choice cascades through
\textsc{Ladder3pt} (wrong kernel generator), the challenge isogeny
(different curve $E_{\text{chl}}$), and every subsequent verification
step.  In our implementation the $j$-invariant of $E_{\text{chl}}$ was
completely wrong until we matched the C reference's canonicalization.

The spec should define a canonical $\Fpii$ square root.  The C reference
uses the convention from Aardal et al.~\cite{ABCDK24}: negate the
output so that the real part (as an integer in $[0, p)$) is even; if the
real part is zero, negate so the imaginary part is even.  This should be
stated in \S8.1 (field arithmetic) as a normative requirement.  A note
of the form
\begin{quote}
  \itshape Throughout this document, ``compute $\sqrt{a}$'' for
  $a \in \Fpii$ means the unique root $r$ such that $r^2 = a$ and
  $\operatorname{Re}(r) \bmod 2 = 0$ (or, if $\operatorname{Re}(r)=0$,
  $\operatorname{Im}(r) \bmod 2 = 0$).
\end{quote}
would suffice.

% =====================================================================
\section{ProductToTheta coordinate convention}
\label{sec:product-to-theta}

\finding{Affects \algo{8.37} (\textsc{ProductToTheta}).}{\high}

The product theta coordinates are described abstractly via level-2 theta
functions.  The mapping from Montgomery projective coordinates $(X:Z)$ to
product theta coordinates is not stated explicitly.  The natural
``theta-like'' representation for Montgomery arithmetic is $(X-Z,\,X+Z)$,
which appears throughout the isogeny literature.  The C reference uses
raw $(X, Z)$ directly.  Using $(X-Z, X+Z)$ produces a completely wrong
change-of-basis matrix~$N$.

\algo{8.37} should state explicitly that the product theta coordinates
are formed as
\[
  (a, b, c, d) = (X_1 X_2,\; X_1 Z_2,\; Z_1 X_2,\; Z_1 Z_2)
\]
from Montgomery projective pairs $(X_1 : Z_1)$, $(X_2 : Z_2)$.

% =====================================================================
\section{``Squared theta'' definition}
\label{sec:squared-theta}

\finding{Affects \algo{8.29}, \algo{8.34}, \algo{8.38}.}{\high}

Several algorithms refer to ``squared theta coordinates'' without a
centralized definition.  In the C reference,
\texttt{to\_squared\_theta(P)} computes component-wise squaring
\emph{followed by a Hadamard transform}: $\tilde{P} =
\mathcal{H}(P^2)$.  The name suggests only squaring.  We omitted the
inner Hadamard in the generic isogeny evaluation, computing
$\mathcal{H}(\mathrm{inv} \cdot P^2)$ instead of the correct
$\mathcal{H}(\mathrm{inv} \cdot \mathcal{H}(P^2))$.

Add a definition early in \S8.5:
\begin{quote}
  \itshape The \textbf{squared theta} of a point $P = (a:b:c:d)$ is
  $\tilde{P} = \mathcal{H}(a^2, b^2, c^2, d^2)$, i.e.\ component-wise
  squaring followed by the Hadamard transform.
\end{quote}

% =====================================================================
\section{\texttt{hadamard\_bool} in the $(2,2)$-chain}
\label{sec:hadamard-bool}

\finding{Affects \algo{8.47}
(\textsc{Isogeny22ChainWithTorsion}).}{\high}

The algorithm describes each generic step uniformly.  The C reference
uses three formula variants controlled by boolean flags:
\begin{itemize}[nosep]
  \item Normal ($i < n{-}2$): codomain Hadamard-transformed;
        eval~$= \mathcal{H}(\mathrm{precomp} \cdot \mathcal{H}(P^2))$.
  \item Penultimate ($i = n{-}2$): codomain in dual form;
        eval~$= \mathrm{precomp} \cdot \mathcal{H}(P^2)$.
  \item Ultimate ($i = n{-}1$): extra input Hadamard, codomain dual;
        eval~$= \mathrm{precomp} \cdot \mathcal{H}(\mathcal{H}(P)^2)$.
\end{itemize}
The splitting step expects dual form.  Using the normal formula for all
steps feeds the splitter a Hadamard-transformed null point with no zero
$U_{i,j}(0)$ entry, causing verification to fail after an otherwise
correct 125-step chain.

A short remark in \algo{8.47} noting that the last two steps omit the
final Hadamard (and the ultimate step pre-Hadamards the input) would
prevent this.

% =====================================================================
\section{Torsion basis slot convention}
\label{sec:basis-swap}

\finding{Affects \algo{2.2}
(\textsc{TorsionBasisFromHint}).}{\medium}

The spec says the algorithm returns $(P, Q, P{-}Q)$.  The C reference
stores $P{-}Q$ in the \texttt{Q} slot and $Q$ in \texttt{PmQ}, so
\textsc{Ladder3pt} computes $P + [m](P{-}Q)$, not $P + [m]Q$.  The swap
ensures the multiplied point avoids $(0,0)$ but is not documented.

State the output convention explicitly, or note that \textsc{Ladder3pt}
is called with the difference in the $Q$ position.

% =====================================================================
\section{Normalization of $(A_{24} : C_{24})$}
\label{sec:normalize-a24}

\finding{Affects \algo{2.2}
(\textsc{TorsionBasisFromHint}).}{\medium}

The C reference normalizes to $((A{+}2)/4 : 1)$ before basis generation.
If the curve arrives from an isogeny codomain with $C_{24} \neq 1$, the
Montgomery ladder produces a different projective representative, which
propagates through \textsc{DifferencePoint} and \textsc{Ladder3pt}.

Note in \algo{2.2} that the curve must be in normalized form before
computing the basis.

% =====================================================================
\section{Gluing requires Jacobian $y$-coordinates}
\label{sec:gluing-jacobian}

\finding{Affects \algo{8.39} (\textsc{GluingEval}).}{\low}

Montgomery $x$-only arithmetic cannot distinguish $P+Q$ from $P-Q$.  The
gluing evaluation needs this distinction.  The C reference lifts to
Jacobian $(x,y,z)$ via Okeya--Sakurai for this step---the only place in
verification that needs $y$-coordinates.  Note this in \algo{8.39}.

% =====================================================================
\section{Jacobian doubling for balanced strategy}
\label{sec:jacobian-doubling}

\finding{Affects \algo{8.47}, Phase~1.}{\low}

The spec does not specify whether Phase~1 doublings use Montgomery or
Jacobian coordinates.  The C reference doubles in Jacobian, then converts
via $\texttt{jac\_to\_xz}{:}\; (x,y,z) \mapsto (x, z^2)$.  The theta
coordinate computations are sensitive to the projective representative.
Note this, or at minimum state that the representative matters.

% =====================================================================
\section{\textsc{DifferencePoint} normalization factor}
\label{sec:difference-normalization}

\finding{Affects \algo{8.12} (\textsc{DifferencePoint}).}{\low}

The C reference multiplies $B_{XX}$, $B_{XZ}$, $B_{ZZ}$ by
$C \cdot \overline{C}^{\,2} \cdot \overline{Z_P}^{\,2} \cdot
\overline{Z_Q}^{\,2}$ (Galois conjugation) before the quadratic solve.
This makes the discriminant a fourth power in~$\Fp$, reducing the
$\Fpii$ sqrt to an $\Fp$ sqrt.  We initially used $(Z_P Z_Q)^2$, which
lacks this property since $\overline{z}^{\,2} \neq z^2$ in general.
State the normalization factor and its purpose in \algo{8.12}.

% =====================================================================

\begin{thebibliography}{9}

\bibitem{sqisign-spec}
{SQIsign team},
\textit{SQIsign: Compact Post-Quantum Signatures from Quaternions and Isogenies},
Specification version 2025-07-07,
\url{https://sqisign.org/spec/sqisign-20250707.pdf}.

\bibitem{ABCDK24}
M.~N.~H.~Aardal, M.~P.~Azimova, E.~L.~Carri\'{o}n, L.~Dillinger, and
M.~J.~Kannwischer,
\textit{Optimized One-Dimensional SQIsign Verification on Intel and Cortex-M4},
Cryptology ePrint Archive, Paper 2024/1563,
\url{https://eprint.iacr.org/2024/1563}.

\bibitem{CHR17}
C.~Costello, H.~Hisil, and B.~Smith,
\textit{Faster Compact Diffie--Hellman: Endomorphisms on the $x$-line},
ASIACRYPT 2017,
\url{https://eprint.iacr.org/2017/518}.

\end{thebibliography}

\clearpage
\appendix

% =====================================================================
\section{Never recompute $P{-}Q$ via \textsc{DifferencePoint}}
\label{sec:pmq-recompute}

\finding{Affects \algo{4.9} (verification), lines 11--26,
and transitively the $(2,2)$-isogeny chain.}{\critical}

A torsion basis $(P,\, Q,\, P{-}Q)$ is produced by
\textsc{TorsionBasisFromHint} (\algo{2.2}).  The third component
$P{-}Q$ is computed once, via \textsc{DifferencePoint}, which
internally takes an $\Fpii$ square root.  Subsequent operations---scaling
(repeated doubling), matrix application (\algo{4.9} line~14), and the
even response isogeny (line~17)---must propagate all three points
together.  If at any point $P{-}Q$ is recomputed from $P$ and $Q$ via
\textsc{DifferencePoint}, the fresh square root may pick the opposite
branch, producing a point that is \emph{affinely distinct} from the
original $P{-}Q$.

Concretely, the two branches of \textsc{DifferencePoint} yield $x(P-Q)$
and $x(P+Q) = x(2P-Q)$; these are different points.  All subsequent
differential additions (the three-point ladder, the biladder for matrix
application, and the $(2,2)$-isogeny chain's \textsc{Ladder3pt} and
\textsc{lift\_basis}) consume $P{-}Q$ as a third input and produce wrong
results if it is silently replaced by $2P{-}Q$.

In our implementation, three separate recomputations caused failures:
\begin{enumerate}[nosep]
  \item After scaling the basis by repeated doubling, we reconstructed
        $P{-}Q$ from the doubled $P$ and $Q$ instead of doubling $P{-}Q$
        alongside them.
  \item The matrix application $M_{\mathrm{chl}} \cdot (P, Q)$ computed
        $P' - Q'$ via \textsc{DifferencePoint}$(P', Q')$ instead of a
        third biladder call $[(a{-}b)]P + [(c{-}d)]Q$.
  \item Before entering the $(2,2)$-chain, we recomputed
        $P{-}Q$ from the post-isogeny images instead of pushing the
        original $P{-}Q$ through the isogeny as a third evaluation point.
\end{enumerate}
Each of these independently caused KAT verification failure.  Fixing all
three was necessary and sufficient to make KAT vector~0 pass.

\textbf{Current spec text.} The spec does not warn against recomputing
$P{-}Q$.  Algorithms that scale or transform a basis mention only $P$
and $Q$; the third point $P{-}Q$ is not discussed.

\medskip\noindent\textbf{Recommendation.}  Add a normative note in \S2.2 or \S8.2:
\begin{quote}
  \itshape Once a basis $(P, Q, P{-}Q)$ has been produced by
  \textsc{TorsionBasisFromHint}, the third point $P{-}Q$ must be
  propagated through every subsequent operation (doubling, matrix
  application, isogeny evaluation) alongside $P$ and~$Q$.
  Recomputing $P{-}Q$ via \textsc{DifferencePoint} is not permitted:
  the square root may select a different branch, yielding an affinely
  distinct point that silently corrupts all downstream differential
  additions.
\end{quote}

% -------------------------------------------------------------------
\section{Change-of-basis matrix application convention}

\finding{Affects \algo{4.8} (\textsc{SetChangeOfBasisMatrix}),
  \algo{4.9} (verification, line~14).}{\high}

The spec does not specify whether the change-of-basis matrix
$\mathbf{M}_\mathrm{chl}$ is applied to a torsion basis $(P, Q)$
by rows or by columns.

\paragraph{What the spec says.}
\algo{4.8} line~6 writes $(P_2, Q_2) \leftarrow \mathbf{M}_1 \cdot (P_2, Q_2)$
without specifying the convention.
\algo{4.9} line~14 writes
$(P_\mathrm{chl}, Q_\mathrm{chl}) \leftarrow \mathbf{M}_\mathrm{chl} \cdot (P_\mathrm{chl}, Q_\mathrm{chl})$,
again without specifying.

\paragraph{What the C reference does.}
\texttt{matrix\_scalar\_application\_even\_basis} in \texttt{verify.c}
applies the matrix by \emph{columns}:
\begin{align*}
  R' &= [a]\,P + [c]\,Q  &&\text{(column 0: \texttt{mat[0][0]}, \texttt{mat[1][0]})} \\
  S' &= [b]\,P + [d]\,Q  &&\text{(column 1: \texttt{mat[0][1]}, \texttt{mat[1][1]})}
\end{align*}
where $\mathbf{M} = \bigl(\begin{smallmatrix} a & b \\ c & d \end{smallmatrix}\bigr)$.
An implementor who applies by rows (the natural reading of
``$\mathbf{M} \cdot (P, Q)^T$'') will produce wrong results.

\paragraph{Impact.}
Applying by rows instead of columns silently produces the wrong
torsion basis, causing the (2,2)-isogeny chain in verification to
fail (typically a panic in the splitting step).

\paragraph{Recommendation.}
State explicitly in \algo{4.8} and \algo{4.9} that the matrix is
applied by columns, or equivalently define the multiplication as
$(P', Q') = (P, Q) \cdot \mathbf{M}^T$.

% =====================================================================
\section{Fixed-precision Hermite Normal Form}
\label{sec:hnf}

\finding{Affects \algo{3.2} (\textsc{HNF}), and transitively every
algorithm that constructs or reduces an ideal lattice --- including
\algo{3.9} (\textsc{RandomEquivalentPrimeIdeal}),
\algo{3.10} (\textsc{RandomIdealGivenNorm}), and
\algo{3.15} (\textsc{FixedDegreeIsogeny}).}{\high}

\algo{3.2} describes the column-style Hermite Normal Form via
extended-GCD elimination.  The pseudocode operates over
$\mathbb{Z}$, i.e., arbitrary-precision integers.

\paragraph{What goes wrong.}
Classical HNF's intermediate entries can grow exponentially in
the rank.  For rank-4 quaternion ideal lattices at NIST-I the
initial column entries can reach ${\sim}2^{770}$~bits (the
$p \cdot g_i$ products from \algo{3.10}).  After a few xgcd
elimination steps, intermediate column entries exceed any
reasonable fixed-width budget.  At our storage width of 1920~bits
($30 \times 64$), the resulting ``HNF'' silently wraps: we
observed basis columns collapse to elements of the ambient
order~$O_0$ (reduced norms of~4 and~${\sim}2^{251}$) rather
than the intended ideal~$I$.  Every subsequent operation ---
\textsc{RandomEquivalentPrimeIdeal}, \textsc{IdealToIsogeny},
the entire response phase --- then operates on a lattice unrelated
to~$I$.

\paragraph{What the C reference does.}
The C reference uses GMP (\texttt{ibz\_t}), so coefficient growth
is absorbed by arbitrary-precision arithmetic.  The classical
algorithm works correctly there.

\paragraph{Fix.}
Use the standard Domich--Kannan--Trotter modular HNF (Cohen,
\emph{A Course in Computational Algebraic Number Theory},
\S2.4.2): given a positive multiple~$D$ of the lattice
determinant~$\det(L)$, reduce every intermediate entry modulo~$D$
after every arithmetic update.  This bounds entries by~$D$ instead
of by the exponential classical bound, at the cost of computing
at a wider intermediate width $W \geq 2 \cdot
\lceil\log_2 D\rceil / 64$.  The result is identical to the
classical HNF in~$\mathbb{Z}$.

For NIST-I commitment ideals, $D = 4 d^4 N^2 p \approx 2^{1282}$
is a compile-time constant.  For response-phase and post-reduce
ideals the modulus varies and is computed per call.

\paragraph{Recommendation.}
Add a remark after \algo{3.2} noting that fixed-precision
implementations must account for coefficient growth, and that the
standard modular HNF technique (reducing modulo a known multiple
of the lattice determinant) is sufficient to keep intermediates
bounded.

\end{document}
